\section{Common ground vectors on disjoint regions}

Let $\mathcal H_{L,q}=(\C^q)^{\otimes L}$, with $q\ge1$ and $L\in\N$.
Support on a site set means that the operator is the identity on its
complement, as in Definition~\ref{def:chain_model}. Empty regions and the
empty chain are allowed.

\begin{lemma}[One set of original minima for every clipping subset]
    \label{thm:disjoint_regional_clipping}
    \lean{Matrix.exists_disjoint_regional_clipping_structure}
    \leanok
    \uses{def:chain_model}
    Let $(R_i)_{i\in I}$ be a finite family of pairwise disjoint site sets,
    and let $A_i$ be supported on $R_i$ with $0\le A_i\le c_i\Id$,
    where $c_i\in\N$. Let $0\le Y\le C\Id$ with $C\in\N$.
    There are real numbers $e_i$ and a unit vector $x$ with
    $A_ix=e_ix$ and $A_i\ge e_i\Id$ for every $i$ such that the following
    assertions hold for every $t\ge0$ and every subset $K\subseteq I$.
    Put
    \begin{align}
        F_i & =\begin{cases}\min(A_i,e_i+t),&i\in K,\\A_i,&i\notin K,\end{cases}
            & X                                                                  & =\sum_i(F_i-e_i\Id),\notag \\
        b   & =\sum_{i\in K}t+\sum_{i\notin K}c_i+C.
    \end{align}
    Then $X\ge0$, $Xx=0$, and $X$ commutes with every $A_j$.
    For every $j$ and $0\le u\le t$,
    $u\mathbf1_{\{\lambda\ge e_j+u\}}(A_j)\le X$.
    Moreover $X+Y\le b\Id$ and $\|X+Y\|\le b$.
    There are a unit vector $z$ and a real number $E$ satisfying
    $(X+Y)z=Ez$, $X+Y\ge E\Id$ and $0\le E\le C$.
    The same $e_i,x$ work for every $t,K$; $E,z$ may depend on $t,K$.
\end{lemma}
\begin{proof}
    \leanok
    Induct over the disjoint regions to obtain a common product of regional
    ground vectors. Their attained energies are the original minima $e_i$.
    Disjoint supports imply commutation. The centered partial-clipping
    inequalities in~\cite[``Centered partial clipping of a finite family'']
    {QICLean2026Blueprint} give the order, high-projection and norm bounds.
    A minimizing unit eigenvector of $X+Y$ exists in finite dimension.
    Testing its minimum against $x$ gives $E\le\langle x,Yx\rangle\le C$.
    This is the disjoint-region construction used in
    \cite[lines 873--886 and 979--1005]{Arad2016RigorousRG}.
\end{proof}
